\subsection{Excluding a five-vertex positive-surplus core}

Throughout this subsection, $D$ is a strongly connected counterexample
which is minimal under deletion of nonempty sets of arcs, and
$n=2\delta+3$. Write
\[
 e(v)=d^+(v)-\delta,\qquad
 g(v)=d^+(v)-d^{++}(v),\qquad b(v)=g(v)-1.
\]
Thus $b(v)\ge0$. For a target $x$, put
\[
 T(x)=V\setminus N^-[x],\quad
 P(x)=\{u:x\in U(u)\},\quad p(x)=|P(x)|,
 \quad q(x)=e(x)+\mu(x).
\]
For $q(x)>0$ define $\sigma(x)=2q(x)-1-p(x)$, and put
$\sigma(x)=0$ when $q(x)=0$. Let
\[
 K=\{x:\sigma(x)>0\},\quad k=|K|,\qquad
 S=\sum_x\sigma(x).
\]
Call $\sigma=1$ targets light and $\sigma\ge2$ targets heavy.
We use the previously established core identity
\begin{equation}\label{eq:core-count-six}
 \binom{k}{2}=S+m+t+A_K+M_{\rm out},
\end{equation}
where $m$ is the number of minimum-degree vertices in $K$,
$t$ counts gap-two vertices of excess at most one,
$A_K$ is the number of arcs induced by $K$, and $M_{\rm out}$
is the number of missing pairs with both endpoints in $V\setminus K$.

\begin{lemma}[Minimum predecessors of heavy targets]\label{lem:heavy-minimum-predecessor}
Every target $s$ with $\sigma(s)\ge2$ has an in-neighbour of
out-degree $\delta$.
\end{lemma}
\begin{proof}
Suppose otherwise. Then $D'=D-s$ has order $n'=2\delta+2$
and all its out-degrees are at least $\delta$: minimum-degree vertices
lose no outgoing arc, and every other vertex loses at most one.
The oriented arc-count bound gives minimum out-degree at most
$\lfloor(n'-1)/2\rfloor=\delta$, so it is exactly $\delta$.
The fixed-target
budget identity in $D'$ gives
\[
 B'+n'=z'-\sum_x\sigma'(x)\le n',\qquad
 B'=\sum_{v\in V(D')}b_{D'}(v),
\]
since the capacity bound gives $\sigma'(x)\ge0$ and $z'\le n'$.
Thus $B'\le0$.

Deleting $s$ can only remove exact second neighbours of the other
vertices. If $E_s(v)$ denotes those lost second neighbours other
than $s$, the exact transfer identity is
\[
 B'=B-b(s)+\bigl(d^{--}(s)-d^-(s)\bigr)
       +\sum_{v\ne s}|E_s(v)|.
\]
The first term $B-b(s)$ and the final sum are nonnegative.
Furthermore
\[
 d^{--}(s)-d^-(s)=2q(s)-2-p(s)=\sigma(s)-1\ge1.
\]
Consequently $B'\ge1$, a contradiction.
\end{proof}

\begin{lemma}[A weighted five-vertex relation]
Let $R$ be an oriented graph with positive integer vertex weights
$w_1,\ldots,w_H$ whose sum is five. If
\[
 \sum_{j\in N_R^+(i)}w_j\ge2\quad\text{for every }i,
\]
then $H=5$, all weights equal one, and $R$ is a regular tournament.
\end{lemma}
\begin{proof}
There is no possibility with $H=1$ or $H=2$. With $H=3$, positive
out-degree at every vertex forces a directed triangle. Each
successor weight would then be at least two, exceeding the total
weight five. With $H=4$, the weights are $2,1,1,1$. The weight-two
vertex must point to two singleton vertices. Neither can point
back to it; each must therefore point to both remaining singleton
vertices, giving a digon. Finally, with $H=5$ every weight is one.
The minimum out-degree two uses all ten possible arcs, and every
out-degree is two.
\end{proof}

\begin{theorem}\label{thm:core-six}
The positive-surplus core satisfies $k\ge6$. In particular,
\[
 {z\ge h+8}.
\]
\end{theorem}
\begin{proof}
An all-deficient graph has minimum outdegree at least two, so the outside tournament used below is nonempty. The preceding core bound gives $k\ge5$. Suppose $k=5$.
The established light-target incidence bound $A_K\ge L$, where
$L$ is the number of light targets, and $S\ge k+H$, where $H$
is the number of heavy targets, make every inequality in
\eqref{eq:core-count-six} an equality. In particular,
\[
 m=t=M_{\rm out}=0,\qquad
 \sigma(x)\in\{1,2\}\ (x\in K),\qquad A_K=L.
\]
Thus every core vertex is nonminimum-degree; every outside vertex
has degree $\delta$ and is a zero-surplus target; and the outside
graph is a tournament on $2\delta-2$ vertices.

 Every light target is nonminimum-degree and, by the previously
 established degree--surplus bounds, has $e=1$, $g=1$, and
 $|U(x)|=1$. A light target with $q=1$, $p=0$ would therefore
 have $\mu=0$ and be adjacent to all four other core vertices.
 This would give $A_K\ge4$, whereas $A_K=L=5-H\le3$ by $H\ge2$.
 Hence every light target is active, with $p>0$.

 The active-light protected-source lemma supplies, for each
 light target $x$, a parent $y\in K$ satisfying
 \[
 y\to x,\qquad N^-(y)\subseteq N^-(x)\setminus\{y\}.
 \]
 Since $q(y)>q(x)$, the chosen parent edges form a forest.
 They have distinct heads and number $L$, so they are all the
 core arcs. A parent $y$ cannot itself be light: its own parent
 $z$ would also satisfy $N^-(z)\subseteq N^-(x)$ and $z\to x$.
 The protected containments would make both $y$ and $z$
 unreachable from $x$ within two steps, contradicting $|U(x)|=1$.
 Thus the forest consists of stars rooted at the heavy targets.
 Let $C_1,\ldots,C_H$ be these star basins, with roots
 $r_1,\ldots,r_H$.
 Put $w_i=|C_i|$, so $\sum_iw_i=5$.

Every minimum-degree vertex has at least two unreachable core
vertices, by the tight-target demand lemma. Protected containment
has two consequences. First, a minimum vertex $u$ with $u\to r_i$
points to every vertex of $C_i$. Second, if $u$ reaches $r_j$
within two steps, it reaches every vertex of $C_j$ within two
steps. Both statements follow by iterating the containment along
a protected tree.

Define a directed relation on the heavy roots by
\[
 i\to j\quad\Longleftrightarrow\quad
 \text{there exists a degree-$\delta$ vertex $u$ with}
 \ u\to r_i\text{ and }r_j\in U(u).
\]
Every heavy root has such a minimum predecessor by the preceding
lemma. The predecessor lies outside $K$, and its at least two
unreachable core vertices lie in basins other than its own.
Therefore
\[
 \sum_{j\in N_R^+(i)}w_j\ge2.
\]
The relation has no digon. Indeed, witnesses $u\to r_i$,
$r_j\in U(u)$ and $v\to r_j$, $r_i\in U(v)$ forbid both
$u\to v$ and $v\to u$, since either arc would give a forbidden
two-step path. But $u$ and $v$ are distinct vertices of the
outside tournament.

The weighted-relation lemma now forces $H=5$, all $C_i$ to be
singletons, and the relation to be a regular tournament. If an
outside vertex $u$ points to $r_i$, its unreachable core set
contains at least two vertices and is contained in $N_R^+(i)$,
which has exactly two vertices. It must equal $N_R^+(i)$.
If $u$ pointed to two roots, their relation out-neighbourhoods
would therefore be equal. Distinct vertices of a tournament
cannot have equal out-neighbourhoods: their mutual arc
distinguishes them. Hence every outside vertex has at most one
out-neighbour in $K$.

Every outside vertex consequently has out-degree at least
$\delta-1$ in the outside tournament. Its average out-degree,
however, is $(2\delta-3)/2=\delta-3/2$, a contradiction.
This excludes $k=5$.

Finally, the established $H\ge2$ gives
\[
 z-h=S\ge k+H\ge8,
\]
as claimed.
\end{proof}
